Data-parallel computing frameworks have become essential for processing large-scale scientific datasets in a variety of domains, including seismic processing.
These frameworks accelerate computation by dividing datasets into independent chunks that execute concurrently across distributed clusters.
However, choosing suitable chunk dimensions remains a fundamental challenge that directly impacts both performance and reliability.
Oversized chunks exhaust worker memory and trigger \ac{OOM} failures that terminate entire jobs, while undersized chunks inflate scheduling overhead, reduce cache locality, and underutilize computational resources.

This paper introduces a novel memory-aware chunking mechanism for Dask, \EBRPD{one of the most}{a} widely adopted distributed computing frameworks in the scientific Python ecosystem.
Our approach predicts peak memory usage from input data shapes and operator characteristics, then automatically derives the largest safe chunk size that respects memory constraints.
The prediction model employs lightweight linear regression trained offline on profiling data collected from representative workloads.
Unlike existing heuristics that apply uniform chunk sizes regardless of operator complexity, our method adapts to the specific memory requirements of each computational operator.

We evaluate our approach through experiments with seismic imaging operators, focusing on the memory-intensive \ac{GST3D} computation.
Experiments on synthetic seismic volumes ranging from 100³ to 400³ voxels demonstrate complete elimination of OOM failures across 768 trials, compared to 31\% failure rate with \dask's auto-chunking and 20\% with evenly-split strategies.
Memory-aware chunking reduces median peak memory consumption by 53\% (from 2.62GB to 1.24GB) at the cost of increased execution time (median 28.1s vs 15.7s).
For large volumes where baseline methods fail entirely, memory-aware chunking enables successful completion with up to 8 workers, demonstrating superior scalability in memory-constrained environments.

The implementation integrates seamlessly with \dask's existing scheduling interface, requiring no modifications to user code or the core framework.
Our results show that predictive, operator-aware chunking provides a practical and robust solution for reliable distributed processing in memory-constrained high-performance computing environments.

\EB{Daniel, o abstract está legal, mas me parece muito longo. Tente reduzir o tamanho. Se necessário, passe parte das informações apresentadas nele para a introdução.}